
Der Befehlssatz der Registermaschine Johnny umfasst die folgenden Befehle:

% \begin{itemize}
%     \item \texttt{TAKE x} kopiert den Wert aus Speicherzelle $x$ in den Akkumulator
%     \item \texttt{ADD x} Akkumulator um Inhalt des Registers $x$ erhöhen
%     \item \texttt{SUB x} Akkumulator um Inhalt des Registers $x$ erniedrigen
%     \item \texttt{SAVE x} Inhalt des Akkumulators im Register $x$ speichern
%     \item \texttt{JMP x} springe zu Register $x$
%     \item \texttt{TST x} fals Register $x$ = 0, springe ins übernächste Register
%     \item \texttt{INC x} Erhöhe Inhalt des Registers $x$ um 1
%     \item \texttt{DEC x} Erniedrige Inhalt des Registers $x$ um 1
%     \item \texttt{NULL x} Setze Inhalt des Registers $x$ auf 0
%     \item \texttt{HLT x} Das Programm hält an.
% \end{itemize}

\begin{verbatim}
    TAKE x : kopiert den Wert aus Speicherzelle x in den Akkumulator
    ADD x : Akkumulator um Inhalt des Registers x erhöhen
    SUB x : Akkumulator um Inhalt des Registers x erniedrigen
    SAVE x : Inhalt des Akkumulators im Register x speichern
    JMP x : springe zu Register x
    TST x : fals Register x = 0, springe ins übernächste Register
    INC x : Erhöhe Inhalt des Registers x um 1
    DEC x : Erniedrige Inhalt des Registers x um 1
    NULL x : Setze Inhalt des Registers x auf 0
    HLT x : Das Programm hält an.
\end{verbatim}

Schreibweise: Wir schreiben [10] wenn wir die Zahl an der Speicherstelle (=Register) 10 meinen.

Johnny arbeitet nur mit nicht-negativen Zahlen: Ist ein Ergebnis kleiner als 0, steht dort 0.


